% ─── Анализ ошибок по периодам (MAE по полугодиям) ─────────────────────────
\begin{frame}{Где модели ошибаются больше: MAE по полугодиям}

  \begin{columns}[T]
    % ─── Левая колонка: таблица ──────────────────────────────────────────
    \column{0.55\textwidth}
    \scriptsize
    \renewcommand{\arraystretch}{1.35}
    \setlength{\tabcolsep}{4pt}
    \begin{tabular}{>{\bfseries}lccccc}
      \toprule
      \textbf{Период} &
      \makecell{\textbf{ARIMA}\\\textbf{static}} &
      \textbf{RF} &
      \textbf{LSTM} &
      \makecell{\textbf{ARIMA}\\\textbf{roll.}} &
      \makecell{\textbf{Вол-}\\\textbf{ть}} \\
      \midrule
      1-е пг 2024$^{*}$ & 40  & 29  & 33 & 13 & 13,1 \\
      \rowcolor{light}
      2-е пг 2024       & 117 & 84  & 25 & 11 & 10,6 \\
      1-е пг 2025       & 268 & 220 & 37 & 16 & 15,4 \\
      \rowcolor{light}
      2-е пг 2025       & 382 & 59  & 62 & 32 & 32,2 \\
      \bottomrule
    \end{tabular}

    \vspace{0.25cm}
    {\tiny $^{*}$с апреля 2024 г.\quad MAE в сумах; «Вол-ть» --- средний модуль дневного изменения курса (сум/день).}

    % ─── Правая колонка: диаграмма ───────────────────────────────────────
    \column{0.45\textwidth}
    \centering
    \begin{tikzpicture}
      \begin{axis}[
        ybar, bar width=6pt,
        width=\textwidth, height=4.6cm,
        enlarge x limits=0.18,
        symbolic x coords={1п24, 2п24, 1п25, 2п25},
        xtick=data,
        x tick label style={font=\tiny},
        ylabel={MAE, сум}, ylabel style={font=\tiny},
        yticklabel style={font=\tiny}, ymin=0, ymax=420,
        legend style={font=\tiny, at={(0.5,1.12)}, anchor=south, legend columns=2, draw=none},
        nodes near coords, nodes near coords style={font=\tiny},
        every node near coord/.append style={anchor=south},
      ]
        \addplot[fill=danger!70, draw=danger]   coordinates {(1п24,40)(2п24,117)(1п25,268)(2п25,382)};
        \addplot[fill=success!70, draw=success] coordinates {(1п24,33)(2п24,25)(1п25,37)(2п25,62)};
        \legend{ARIMA static, LSTM}
      \end{axis}
    \end{tikzpicture}

    {\tiny\color{mygray}ARIMA static растёт с горизонтом, LSTM держит точность.}
  \end{columns}

  \vspace{0.2cm}
  \begin{tcolorbox}[infobox]
    \small \textbf{Вывод:} ошибки концентрируются в периоды повышенной волатильности. Модели со свежими данными (ARIMA rolling, RF) и нейросеть LSTM держат точность, тогда как статическая ARIMA деградирует с ростом горизонта прогноза (40\,$\to$\,382 сум).
  \end{tcolorbox}
\end{frame}